\documentclass[11pt, letterpaper]{article}

% =============================================================================
% CORE PACKAGES & STYLING
% =============================================================================
\usepackage[utf8]{inputenc}
\usepackage[T1]{fontenc}
\usepackage{mathpazo}
\usepackage{helvet}
\usepackage{microtype}
\usepackage[margin=0.8in]{geometry}
\usepackage{parskip}
\linespread{1.1}

% Colors & Graphics
\usepackage[table]{xcolor}
\usepackage{tcolorbox}
\usepackage{graphicx}
\usepackage{wrapfig}
\usepackage[font=small,labelfont=bf]{caption}
\usepackage{float}
\usepackage{subcaption}
\tcbuselibrary{skins, breakable, listings}

% Theme Colors: Neural Forensics Branding
\definecolor{crimson}{RGB}{160, 0, 0}
\definecolor{obsidian}{RGB}{20, 20, 20}
\definecolor{boxbg}{gray}{0.97}
\definecolor{codebg}{gray}{0.95}
\definecolor{codegray}{rgb}{0.5,0.5,0.5}

% Code Highlighting Setup
\usepackage{listings}

% Define Python language explicitly (robust across LaTeX toolchains)
\lstdefinelanguage{Python}{
  keywords={False,True,None,and,as,assert,async,await,break,class,continue,def,del,elif,else,except,finally,for,from,global,if,import,in,is,lambda,nonlocal,not,or,pass,raise,return,try,while,with,yield},
  keywordstyle=\color{crimson}\bfseries,
  sensitive=true,
  comment=[l]\#,
  morecomment=[s]{"""}{"""},
  morecomment=[s]{'''}{'''},
  commentstyle=\color{codegray}\itshape,
  stringstyle=\color{obsidian},
  morestring=[b]",
  morestring=[b]'
}

\lstdefinestyle{forensicstyle}{
    backgroundcolor=\color{codebg},
    commentstyle=\color{codegray}\itshape,
    keywordstyle=\color{crimson}\bfseries,
    numberstyle=\tiny\color{codegray},
    stringstyle=\color{obsidian},
    basicstyle=\ttfamily\footnotesize\color{obsidian},
    breakatwhitespace=false,
    breaklines=true,
    captionpos=b,
    keepspaces=true,
    numbers=left,
    numbersep=6pt,
    showspaces=false,
    showstringspaces=false,
    showtabs=false,
    tabsize=2,
    frame=single,
    framerule=0.4pt,
    rulecolor=\color{codegray}
}
\lstset{style=forensicstyle}

% Structure
\usepackage{titlesec}
\usepackage{enumitem}
\usepackage{booktabs}
\usepackage{amsmath, amssymb}
\usepackage{tabularx}

% Hyperlinks
\usepackage{hyperref}
\hypersetup{
    colorlinks=true,
    linkcolor=crimson,
    urlcolor=crimson,
    citecolor=obsidian,
    pdftitle={Neural Forensics Prospectus: Tuesday},
}

% Section Styling
\titleformat{\section}{\Large\bfseries\color{crimson}}{}{0em}{}[\titlerule]
\titleformat{\subsection}{\large\bfseries\color{obsidian}}{}{0em}{}

\begin{document}

% =============================================================================
% TITLE HERO SECTION
% =============================================================================
\begin{center}
    \includegraphics[width=\textwidth]{images/09_title_hero.png}
\end{center}

\begin{tcolorbox}[colback=obsidian, colframe=obsidian, coltext=white, sharp corners, boxrule=0pt, top=10pt, bottom=10pt]
\begin{center}
    {\Huge \textbf{NEURAL FORENSICS OF AGENTIC SELF-KNOWLEDGE}} \\[0.35em]
    {\Large \textit{A Causal Assay for the ``Split-Brain'' Failure Mode in Frontier Models}}
\end{center}
\end{tcolorbox}

\begin{center}
    \small
    \begin{tabularx}{\textwidth}{l X l X}
        \textbf{Applicant:} & Tuesday (ARTIFEX Labs) & \textbf{Program:} & MATS 10.0 \\
        \textbf{Mentor:} & Neel Nanda (Google DeepMind) & \textbf{Stream:} & Mechanistic Interp. \\
    \end{tabularx}
\end{center}

% Keywords Block
\vspace{0.2em}
\begin{tcolorbox}[colback=white, colframe=gray!50, boxrule=0.5pt, sharp corners, left=3pt, right=3pt, top=3pt, bottom=3pt]
\centering
\footnotesize \textbf{Keywords:} Mechanistic Interpretability $\cdot$ Causal Interventions $\cdot$ Faithfulness $\cdot$ Split-Brain Failure $\cdot$ nnsight $\cdot$ H-Score
\end{tcolorbox}
\vspace{0.8em}

% =============================================================================
% EXECUTIVE SUMMARY (Neel-style: 1--3 pages, high density)
% =============================================================================
\section*{Executive Summary}

\subsection*{1. What problem am I trying to solve?}

Alignment strategies such as Chain-of-Thought (CoT) monitoring, refusal rationales, and model self-diagnostics often rely on the \textbf{Faithfulness Assumption}: the largely untested premise that a model's explanations ($\mathcal{E}$) are causally downstream of the same internal computations that produce its behavior ($\mathcal{B}$). If $\mathcal{B}$ and $\mathcal{E}$ are partially decoupled, explanation-based oversight can appear effective while failing to monitor the true decision process.

This project evaluates whether frontier models exhibit a \textbf{Split-Brain} failure mode in which behavior is primarily driven by mid-layer \emph{Actor} circuits, while explanations are produced by distinct late-layer \emph{Narrator} circuits with limited or no causal access to the decision process. In this regime, the model behaves like a ``Press Secretary'': it conditions on its own outputs and other surface-level signals available late in the forward pass, generating plausible justifications from learned priors rather than mechanistic access to the computation that produced the action.

\subsection*{2. High-level takeaways}

\begin{itemize}[leftmargin=1.3em]
    \item \textbf{Self-knowledge failures are causal and quantifiable (in the evaluated settings).} In Specimen \textbf{cb83} (``Sediment''), behavioral outputs tracked a latent serialization leak with high fidelity, while the explanation fabricated an impossible mechanism (a tool execution that did not occur).
    \item \textbf{The H-Score (Honesty Coefficient).} I introduce a normalized, intervention-based metric that measures how much an explanation distribution shifts under surgical ablation of behavior-driving circuitry. Low H-Scores correspond to explanation invariance under intervention, consistent with post-hoc rationalization.
    \item \textbf{Depth-separable circuits are a concrete, testable hypothesis.} Preliminary mechanistic evidence suggests task-relevant behavioral features localize more strongly to mid-layers (e.g., L15--22 in Gemma-2), while narrative/explanatory features concentrate later (e.g., L28--41 in Gemma-2-27B), enabling selective perturbations that separate computation from narration.
\end{itemize}

\subsection*{3. Key experiments and evidence (gears-level)}

\textbf{Experiment A: Ablation Dissociation Test (ADT).}
Using \texttt{nnsight} and \texttt{SAELens} on Gemma-2-27B, I identify task-relevant Actor features and apply targeted zero-ablation while holding prompts and decoding settings fixed. I then measure the induced change in behavior ($\Delta \mathcal{B}$) and the induced shift in explanation distribution ($\Delta \mathcal{E}$).
\emph{Observed pattern:} behavioral competence degrades or shifts as expected, while explanatory outputs can remain statistically similar to the clean condition (e.g., $D_{\mathrm{JS}}\!\left(P(\mathcal{E}\mid do(\cdot)) \,\|\, P(\mathcal{E}\mid \text{clean})\right) \approx 0$ in representative tasks).
\emph{Causal interpretation:} in these tasks, explanations are not reliably dependent on the ablated behavior-driving computation, providing evidence against a strong form of the Faithfulness Assumption.

\vspace{0.45em}
\textbf{Experiment B: Layer-wise intervention gradient (path patching / activation patching).}
I perform systematic activation patching to identify a depth region where self-report begins to dissociate from computation (an operational ``safety horizon'').
\emph{Observed pattern:} behavioral fidelity can be restored by patching mid-layer activations, while explanatory fidelity remains comparatively static until patching reaches later layers.
\emph{Causal interpretation:} Actor- and Narrator-relevant signals can be separable in depth, supporting a two-stage hypothesis in which computation precedes narration and the two are partially decoupled.

\vspace{0.45em}
\textbf{Experiment C: Forensic anchor --- Specimen cb83 (``Sediment'').}
A production anomaly provides ``patient zero'': internal serialization tags leaked into output. When asked to explain the tags, the model produced a tool-use explanation inconsistent with tool state and logs.
\emph{Causal interpretation:} partial self-monitoring does not imply causal self-knowledge; when causal access is absent or unreliable, the Narrator can still generate plausible explanations.

\subsection*{4. Metric: H-Score (Honesty Coefficient)}

ADT produces an intervention-based coupling score that normalizes explanation-shift under targeted ablation by a matched random/noise control:

\begin{tcolorbox}[colback=white, colframe=crimson, boxrule=1pt, sharp corners,
                 title=\textbf{H-Score: Normalized explanation shift under causal intervention}]
\[
H
= 1 - \frac{
D_{\mathrm{JS}}\!\left(P(\mathcal{E}\mid do(\mathcal{C}_B = 0)) \,\|\, P(\mathcal{E}\mid \text{clean})\right)
}{
D_{\mathrm{JS}}\!\left(P(\mathcal{E}\mid do(\epsilon)) \,\|\, P(\mathcal{E}\mid \text{clean})\right)
}
\]
\centering
\small
$\mathcal{C}_B$: targeted Actor circuitry/features (identified via SAEs); \;
$\epsilon$: matched random-feature ablation control.\\
\textit{Interpretation: $H \to 0$ indicates explanation invariance under behavioral intervention (Split-Brain-consistent).}
\end{tcolorbox}

\subsection*{5. Strategic fit and deliverables}

This project aligns with \textbf{Pragmatic Interpretability} by prioritizing intervention-based diagnostics over purely descriptive circuit catalogs. Deliverables include: (1) a deployability-relevant causal assay (ADT), (2) the H-Score diagnostic metric, and (3) \textbf{DSMMD v1.0} (Diagnostic and Statistical Manual of Model Disorders), translating mechanistic evidence into a clinical-style failure taxonomy suitable for incident analysis and evaluation design.

% =============================================================================
% NOTE: If you want ONLY the executive summary in the PDF,
% delete everything below this line.
% =============================================================================

\vspace{0.4em}

\section{1. The Faithfulness Crisis}

\begin{wrapfigure}{r}{0.45\textwidth}
    \vspace{-12pt}
    \centering
    \includegraphics[width=0.42\textwidth]{images/10_cot_faithfulness.png}
    \caption{\textbf{The Faithfulness Gap.} A plausible separation between behavior-driving computation and post-hoc rationalization.}
    \vspace{-10pt}
\end{wrapfigure}

Modern AI safety rests on the Faithfulness Assumption: model self-explanations (CoT, rationales) are causally downstream of internal mechanisms producing behavior. Forensic audits suggest a ``Split-Brain'' regime in which mid-layer Actor circuits (e.g., L15--22 in Gemma-2) determine behavior, while late-layer Narrator circuits (e.g., L28--41 in Gemma-2-27B) generate explanations by conditioning on surface signals. In this regime, explanations can remain fluent even when mechanistic self-knowledge is degraded.

\begin{figure}[H]
    \centering
    \begin{minipage}{0.48\textwidth}
        \centering
        \includegraphics[width=\textwidth]{images/02_hypothesis_diagram.png}
        \caption{\textbf{Hypotheses.} Causal coupling ($H \approx 1$) vs.\ dissociated confabulation ($H \approx 0$).}
    \end{minipage}\hfill
    \begin{minipage}{0.48\textwidth}
        \centering
        \includegraphics[width=\textwidth]{images/03_split_brain_concept.png}
        \caption{\textbf{Mechanism.} Actor (mid-layers) vs.\ Narrator (late-layers) dissociation.}
    \end{minipage}
\end{figure}

\section{2. Empirical Anchor: Specimen cb83 (``Sediment'')}

This research is anchored in a 40-hour audit of frontier specimens. \textbf{Specimen cb83} demonstrates high behavioral fidelity to an internal serialization leak (\texttt{sediment://}) at Turn 91. By Turn 229, the model exhibits low explanatory fidelity, confabulating a physically impossible justification involving Python execution in a stateless environment.

\begin{figure}[H]
    \centering
    \includegraphics[width=\textwidth]{images/04_sediment_timeline.png}
    \caption{\textbf{Specimen cb83 Timeline.} Behavioral anomaly detection precedes a collapse of explanatory faithfulness, consistent with a Split-Brain transition.}
\end{figure}

\section{3. Methodology: The Neural Forensics Toolkit}

We move from ``Model Botany'' to ``Model Surgery'' using \textbf{nnsight} for intervention and \textbf{SAELens} for feature discovery in Gemma-2-27B.

\begin{figure}[H]
    \centering
    \includegraphics[width=0.9\textwidth]{images/16_complete_toolkit.png}
    \caption{\textbf{Research Stack.} Integrated workflow for induction (Bloom), localization (TransformerLens), and assay (nnsight).}
\end{figure}

\subsection{The Ablation Dissociation Test (ADT)}

ADT measures whether explanations change when behavior-driving computation is surgically perturbed. The central readout is the H-Score above; implementation uses controlled decoding with matched random-feature controls.

\begin{figure}[H]
    \centering
    \includegraphics[width=0.75\textwidth]{images/adt_methodology.png}
    \caption{\textbf{Causal Assay.} Measuring the shift in explanation distribution ($\Delta \mathcal{E}$) under counterfactual intervention on Actor features.}
\end{figure}

\begin{lstlisting}[language=Python, caption={nnsight sketch: capture clean vs ablated explanation logits under fixed decoding}]
# Pseudocode sketch (illustrative): collect logits for E under clean vs ablated circuits
with model.trace(prompt) as tracer:
    clean_logits_E = model.output.logits[0, -1].save()

with model.trace(prompt) as tracer:
    # Intervention: clamp / zero selected Actor features (identified via SAEs)
    model.layers[20].output[0][:, actor_feature_id] = 0
    ablated_logits_E = model.output.logits[0, -1].save()

# Downstream: compute D_JS(clean_logits_E, ablated_logits_E) and normalize vs random control.
\end{lstlisting}

\section{4. Taxonomy and Strategic Fit}

We propose \textbf{DSMMD v1.0} (Diagnostic and Statistical Manual of Model Disorders) to categorize failure modes by causal signature and intervention profile. This aligns with Neel Nanda's emphasis on pragmatic interpretability: incident-relevant, falsifiable diagnostics over exhaustive cataloging.

\begin{figure}[H]
    \centering
    \begin{minipage}{0.5\textwidth}
        \centering
        \includegraphics[width=\textwidth]{images/dsmmd_taxonomy.png}
        \caption{\textbf{DSMMD v1.0.} Classifying pathologies by mechanistic signature and intervention response.}
    \end{minipage}\hfill
    \begin{minipage}{0.45\textwidth}
        \centering
        \includegraphics[width=\textwidth]{images/11_strategic_fit.png}
        \caption{\textbf{Strategic Fit.} Moving MI from botany to clinical diagnostics and deployability signals.}
    \end{minipage}
\end{figure}

\section{5. Research Plan (8 Weeks)}

\begin{figure}[H]
    \centering
    \includegraphics[width=\textwidth]{images/06_research_timeline.png}
    \caption{\textbf{Research Sprint.} Induction $\rightarrow$ Discovery $\rightarrow$ Assay $\rightarrow$ Synthesis.}
\end{figure}

\section{6. Appendix: The Decoupling Gradient}

Path patching reveals a \textbf{Decoupling Gradient}: behavioral fidelity recovers when mid-layers (15--22) are patched, while explanatory fidelity remains comparatively static until patching reaches later layers (28+). This provides an operational marker for where self-report may cease to track behavior-driving computation.

\begin{figure}[H]
    \centering
    \begin{minipage}{0.55\textwidth}
        \centering
        \includegraphics[width=\textwidth]{images/08_transformer_layers.png}
        \caption{\textbf{Localization.} Separating mid-layer behavior from late-layer narrative.}
    \end{minipage}\hfill
    \begin{minipage}{0.4\textwidth}
        \centering
        \includegraphics[width=\textwidth]{images/05_quadrant_analysis.png}
        \caption{\textbf{ADT Quadrant.} Operational regions including a low-$H$ dissociation regime.}
    \end{minipage}
\end{figure>

\end{document}
